\section{Nilpotent covectors and transverse Hecke modifications}
\label{geom:section}

The sheaf-theoretic detection argument needs two geometric inputs: a
dimension bound for the nilpotent cotangent cone and a transverse Hecke
modification for every covector outside it.  We prove these for all the
level stacks of Section~\ref{found:level-stacks}.  Throughout this
section the ground field $\Bbbk$ is algebraically closed of characteristic zero
or of characteristic $p>n$; the product and quotient stacks are needed
only in characteristic zero.  Put $\mathfrak g=\Lie G$, and write
$\mathfrak m=\Lie M$ for a Levi subgroup $M$.

\subsection{The matrix hypotheses and the cotangent cone}

The trace form $\kappa(a,b)=\operatorname{tr}(ab)$ is nondegenerate on
$\mathfrak{sl}_n$, since $n$ is invertible.  We record the other Lie
algebra facts that underlie \cite[\S3]{CT}, with the characteristic
hypotheses checked for our group.  A Levi has block sizes
$d_1,\ldots,d_r$ with sum $n$.  Its Lie center consists of block scalars
$(c_i)$ with $\sum_i d_ic_i=0$.  The trace form there is the restriction
of $\sum_i d_ic_ic'_i$.  Each $d_i$ and $n$ is invertible, so this
restriction is nondegenerate: its orthogonal complement in the block
scalar space is the line $(1,\ldots,1)$, whose norm is $n$.

Chevalley restriction for such a Levi follows from the characteristic
polynomials of the blocks.  The Levi Weyl group has order
$\prod_i d_i!$ invertible in the field, so restriction to the trace-zero
hyperplane commutes with taking invariants.  Symmetric polynomials in
the block eigenvalues give all invariants there; injectivity follows
from the density of matrices with distinct eigenvalues within the
blocks.  Conjugators may be chosen with determinant one by adjusting
a diagonal centralizer element.  The scheme centralizer of a semisimple
matrix is the corresponding block Levi.  Finally, matrix nilpotence
agrees with the zero fiber of these invariants, and every nilpotent
matrix over \emph{any} field extension has a rational full flag on
which it is strictly triangular.  Such a flag is obtained successively
from nonzero vectors in the kernel and the induced nilpotent operator
on the quotient.  These verify the four group hypotheses used in
\cite[\S\S1 and~3]{CT}, without extending the field to find a nilpotent
flag.

Let $P$ be a bundle with reductions $H_i$ at $x_i$, where each $H_i$
is a Borel or its unipotent radical.  Its deformation bundle is
\[
 \mathcal E_H=\ker\left(\ad(P)\longrightarrow
              \bigoplus_i\ad(P)_{x_i}/\mathfrak h_i\right),
 \qquad \mathfrak h_i=\Lie H_i.
\]
Infinitesimal automorphisms and deformations are $H^0(\mathcal E_H)$
and $H^1(\mathcal E_H)$.  Serre duality identifies degree-zero
cotangent vectors with Higgs fields
\begin{equation}\label{geom:residues}
 \phi\in H^0(X,\ad(P)\otimes\omega_X(D)),\qquad
            \Res_{x_i}\phi\in\mathfrak h_i^\perp.
\end{equation}
The trace form gives
$\mathfrak b_i^\perp=\Lie N_{B_i}$ and
$(\Lie N_{B_i})^\perp=\mathfrak b_i$.  Thus ordinary flags impose
nilradical residues, whereas enhanced flags impose Borel-valued
residues.  For unlevelled bundles the fields are regular.

On $\cB_{12}$ a cotangent vector pulls back to a pair of enhanced-level
fields that annihilates the diagonal torus action.  Pairing a covector
with an infinitesimal change of enhancement is pairing with the toral
part of its residue.  The two branch contributions are added with the
signs fixed by the identifications $\tau_i$ in
\eqref{found:quotient}.

\begin{definition}\label{geom:cone}
For any of these stacks $\cB$, let $\Lambda\subset T^*\cB$ be the cone
of Higgs fields generically nilpotent on each component curve in its
presentation.  For a smooth chart $f:S\to\cB$, write
$\Lambda_S=f^\circ\Lambda$.
\end{definition}

The cone is closed, since the nonleading coefficients of the
characteristic polynomial are global sections whose vanishing is
exactly generic nilpotence.  Its enhanced-level residues have zero
toral part.  Indeed, writing $\phi=b(t)\,dt/t$ at a marked point,
the characteristic coefficients of $b(t)$ vanish and hence those of
$b(0)$ vanish.  A nilpotent upper triangular matrix has zero diagonal.
Consequently the enhanced cone is precisely the smooth cotangent
pullback of the ordinary cone under $q$.  Likewise, the cone on
$\cB_{12}$ is the smooth cotangent pullback from
$\cA_1\times\cA_2$.  All toral residues vanish already, so the
quotient imposes no additional condition on this cone.

\begin{proposition}\label{geom:dimension}
For a smooth finite-type chart $S\to\cB$ of pure dimension $d$,
\[
                            \dim\Lambda_S\le d.
\]
In characteristic zero the cone is isotropic: the symplectic form
vanishes on the smooth locus of every reduced subvariety of
$\Lambda_S$.
\end{proposition}

\begin{proof}
We adapt the rational-reduction proof of \cite[Proposition~3.1]{CT}
to all the marked flags.  This is the parabolic form of Ginzburg's
isotropy argument \cite[Lemma~7]{GinzburgNilpotent}; compare the
separable-lifting criterion and its bundle-stack application in
\cite[Appendix~D.1.5 and~D.1.8]{AGKRRV}.  For $\mathrm{SL}_n$ the
required flag is already rational.  First take ordinary flags, and let $Z$ be a
reduced irreducible component of $\Lambda_S$, with function field $F$.
Its generic point supplies $(P,(\beta_i),\phi)$ on $X_F$, together
with its lift to $S$.  Over $F(X)$, choose a rational full flag in the
underlying rank-$n$ vector space that makes the nilpotent field strictly
triangular.  Properness of the associated flag bundle extends this
Borel reduction to the whole smooth curve $X_F$.  Denote it by $Q$.
Its nilradical bundle is a subbundle of $\ad(P)$, so generic
containment extends to
\[
 \phi\in H^0\bigl(X_F,
       (Q\times^B\Lie N_B)\otimes\omega_{X_F}(D)\bigr).
\]
No separable or inseparable extension of $F$ has been used.

For each $i$, fix the relative position of $Q_{x_i}$ and $\beta_i$.
Let $\mathcal Y$ classify a global Borel reduction and flags with
these prescribed relative positions.  This is a smooth stack:
$\Bun_B$ is smooth by the vanishing of curve $H^2$, and its added
fibers are products of Borel orbits in $G/B$.  These orbits are smooth,
since intersections of two Borels have their smooth torus-and-root-group
description.  Its deformation bundle consists of sections of
$\ad(Q)$ whose values at $x_i$ lie in
$\Lie Q_{x_i}\cap\mathfrak b_{\beta_i}$.

The pullback of $\phi$ as a covector on $\mathcal Y$ is zero.  Its
pairing with this deformation bundle has no pole at any $x_i$ by the
original residue conditions.  Generically the pairing is zero because
the nilradical annihilates the Borel Lie algebra under the trace form.
It is therefore zero everywhere, and Serre duality gives the asserted
vanishing on deformations.

Spread the reduction and the fixed positions to a dense open of $Z$,
then restrict to its smooth locus.  The canonical one-form on $T^*S$
vanishes on this open: it evaluates $\phi$ on the induced bundle
deformation, which factors through $\mathcal Y$.  Directions relative
to $\cB$ are also annihilated, since the chart covector comes from
$T^*\cB$.  Its exterior differential, the symplectic form, vanishes
there as well.  An isotropic subspace in a $2d$-dimensional symplectic
space has dimension at most $d$, proving the bound.  The unlevelled
case is the same proof with the conditions at the $x_i$ omitted.

Products preserve the bound.  A smooth cotangent pullback adds the
relative dimension of the smooth map to the dimension of the cone;
its canonical one-form is the pullback of the original one.  The cone
descriptions above therefore give the enhanced and quotient cases.
In characteristic zero the rational-flag argument may start with any
reduced irreducible subvariety of $\Lambda_S$, giving the stated
isotropy on its smooth locus.
\end{proof}

\subsection{A detecting cocharacter}

A nonnilpotent field can be detected by moving a modification defined
by an integral cocharacter.  The next matrix argument is the
$\mathrm{SL}_n$ form of \cite[Lemma~3.2]{CT}.

\begin{lemma}\label{geom:cocharacter}
If $a(t)\in\mathfrak{sl}_n[[t]]$ and $a(0)$ is not nilpotent, then,
after a change of formal $\mathrm{SL}_n$-frame, there are a block Levi
$M$ and a cocharacter $\lambda:\mathbb G_m\to Z(M)$ such that
\begin{equation}\label{geom:normal-form}
 a(t)\in\mathfrak m[[t]],\qquad
 \ad(a(0))|_{\mathfrak g/\mathfrak m}\text{ is invertible},\qquad
                    \kappa(a(0),d\lambda)\ne0.
\end{equation}
\end{lemma}

\begin{proof}
Write $a(0)=s+v$ for its commuting semisimple and nilpotent parts.
Both have trace zero, and $s\ne0$.  Diagonalize $s$ and let $M$ be
its block centralizer.  On the off-block matrices $\ad(s)$ is
invertible, whereas $\ad(v)$ is nilpotent and commutes with it.
Their sum is therefore invertible there.

For block sizes $d_i$, the differentials of central cocharacters span
$\mathfrak z(\mathfrak m)=\{(c_i):\sum d_ic_i=0\}$.  Explicitly,
the integer vectors having entries $d_j,-d_i$ in positions $i,j$
and zero elsewhere lie in the cocharacter lattice, and span this
hyperplane over the ground field because the $d_i$ are invertible.
The trace form on the center is nondegenerate, as checked above, and
$v$ is orthogonal to all block scalars.  Some integral central
cocharacter thus has $\kappa(a(0),d\lambda)=\kappa(s,d\lambda)\ne0$.

Finally remove the off-block coefficient of $t^r$ inductively for
$r\ge1$.  Conjugation by an element congruent to $1+t^rY$ modulo
$t^{r+1}$ changes that coefficient by $[Y,a(0)]$.  The invertibility
on off-block matrices supplies $Y$, and smoothness of $\mathrm{SL}_n$
supplies a group element with this first-order value.  The successive
changes converge in $\mathrm{SL}_n[[t]]$, giving the first assertion
of \eqref{geom:normal-form} without changing the constant term.
\end{proof}

\subsection{The two nondegenerate zeros}

Let $H$ be an exact-position Hecke correspondence, with maps
\[
 p:H\to\cB,\quad o:H\to\cB,\quad l:H\to U,
 \qquad \widetilde p=(p,l),\quad q_H=(o,l).
\]
The symbols $p,o$ mean input and output.  On a product or quotient,
$U$ here denotes $U_1\sqcup U_2$.  For relative position $\lambda^+$,
both $\widetilde p$ and $q_H$ are smooth of relative dimension
$m=\langle2\rho_G,\lambda^+\rangle$.

\begin{proposition}\label{geom:transverse}
For a geometric covector $A\in T_b^*\cB\setminus\Lambda$, there are
an allowed leg $y$, such an $H$, a point $h$ with $p(h)=b$ and $l(h)=y$,
a covector $A'\in T^*_{o(h)}\cB$, and $0\ne\xi\in T_y^*U$ with
\begin{equation}\label{geom:identity}
                         p^*A=q_H^*(A',\xi)\quad\text{at }h.
\end{equation}
The following sections have nondegenerate zeros at $h$:
\begin{enumerate}
\item on $H_{\mathrm{in}}=\widetilde p^{-1}(b,y)$, the restriction
of $p^*A$ to the $q_H$-vertical tangent bundle;
\item on $H_{\mathrm{out}}=q_H^{-1}(o(h),y)$, the restriction
of $o^*A'$ to the $\widetilde p$-vertical tangent bundle.
\end{enumerate}
Each restriction is a section of the dual of the indicated vertical
tangent bundle.  The fibers include the identifications of the fixed
bundles; these fibers and the vertical tangent bundles have dimension
and rank $m$, respectively.  The point $y$
can be chosen in any prescribed nonempty open of a component on which
$A$ is not generically nilpotent.
\end{proposition}

\begin{proof}
We give the local calculation of \cite[Proposition~3.3]{CT} and explain
its compatibility with all the levels.  Choose $y$ away from every
marking where the field representing $A$ is nonnilpotent.  Write it
as $a(t)\,dt$ in a formal frame.  Apply
Lemma~\ref{geom:cocharacter} and modify by $\lambda(t)$, with adjoint
lattices
\[
 L=\mathfrak g[[t]],\qquad L'=\operatorname{Ad}(\lambda(t))L,
 \qquad J=L\cap L'.
\]
As $a(t)\in\mathfrak m[[t]]$ and $\lambda$ centralizes $M$, this field
is regular in both lattices.  It extends to an output field $A'$ with
the same residues at all marked points.

The two fixed-side tangent spaces and their residue pairing are
\begin{equation}\label{geom:tangents}
 T_hH_{\mathrm{in}}=L/J,\qquad T_hH_{\mathrm{out}}=L'/J,
 \qquad (D,C)\longmapsto\Res\kappa(D,C)\,dt.
\end{equation}
Indeed input disc gauge transformations move $L'$, output disc gauge
transformations move $L$, and their common stabilizer has Lie algebra
$J$.  For the integer weight decomposition
$\mathfrak g=\bigoplus_j\mathfrak g_j$ of $\lambda$, these quotients are
\begin{equation}\label{geom:truncations}
 L/J=\bigoplus_{j>0}\mathfrak g_j\otimes \Bbbk[[t]]/(t^j),\qquad
 L'/J=\bigoplus_{j<0}\mathfrak g_j\otimes t^j\Bbbk[[t]]/\Bbbk[[t]].
\end{equation}
Opposite weights pair perfectly, and $t^i$ pairs with $t^{-1-i}$.
Thus \eqref{geom:tangents} is a perfect pairing.

Holding output and leg fixed, a vector $C\in L'$ changes the input
by its principal part.  Its pairing with $p^*A$ is
$\Res\kappa(a(t),C)\,dt=0$, because $a(t),C\in L'$ and $L'$ is its
own annihilator under residue.  The smooth cotangent sequence for
$q_H$ now says that $p^*A$ descends to output bundle times leg.  Its
bundle component is $A'$: principal-part deformations supported away
from $y$ and the markings test this assertion where the bundles are
identified.  They span the deformation space, since adding a
sufficiently large pole divisor there kills $H^1$ of the deformation
bundle.  For the leg component replace $t$ by $t-z$ while holding
output fixed.  The logarithmic derivative of the transition
$\lambda(t-z)$ is $-d\lambda/t$.  Serre duality therefore gives
\[
                       \xi=\pm\kappa(a(0),d\lambda)\,dz\ne0,
\]
where the common sign depends only on the gluing convention.  This
proves \eqref{geom:identity}.  It also shows that the second section
vanishes at $h$, since on $\widetilde p$-vertical vectors both
$p^*A$ and the leg differential vanish.

To compute the derivative of the first section, move the output
lattice in the input fiber in direction $D\in L/J$ and transport a
$q_H$-vertical test vector $C\in L'/J$ with it.  Differentiation gives,
up to sign,
\begin{equation}\label{geom:hessian}
 \Res\kappa(a(t),[D,C])\,dt
                    =\Res\kappa([a(t),D],C)\,dt.
\end{equation}
This is independent of the lifts: $\ad(a(t))$ preserves $L,L'$ and
$J$.  It preserves every $\lambda$-weight summand, and on each nonzero
weight its constant term is invertible by \eqref{geom:normal-form}.
It is consequently invertible on the truncated modules in
\eqref{geom:truncations}.  The perfect pairing
\eqref{geom:tangents} shows that \eqref{geom:hessian} is perfect.
It is exactly the derivative of the section at its zero.  For the
second section, hold the output fixed and move $L$ by $C\in L'/J$.
The calculation exchanges $C$ and $D$, giving the transpose of the
same perfect pairing, up to sign.  Both zeros are therefore
nondegenerate and isolated.  Only integer truncation lengths occur;
no cocharacter weight is inverted in the ground field.

Every calculation was made at $y$.  Away from $y$ all the flags and
their residue conditions are unchanged, so the proof applies with any
number of ordinary or enhanced levels.  On a product it acts in one
factor.  On the quotient it descends from that calculation: unchanged
residues preserve annihilation of the diagonal torus, and the Hecke
correspondence is obtained by the smooth quotient base change of
Section~\ref{found:level-stacks}.
\end{proof}

The two zeros serve different purposes below.  The zero on the output
fiber isolates a contribution to the proper Hecke push, while the zero
on the input fiber gives coordinates for the split quadratic equation
that recovers the original vanishing-cycle test.
